% ============================================================
%  CHAPTER 8 — Conclusion
% ============================================================

\chapter{Conclusion}
\label{ch:conclusion}

% ============================================================
\section{Summary}
% ============================================================

This project designed, implemented, and validated a complete ESP32-based embedded
acquisition node for a hospital room and patient monitoring system.
The firmware acquires temperature, humidity, heart rate, and blood oxygen saturation
from two sensor peripherals, applies local signal processing to derive clinically
meaningful values, and transmits authenticated records to Firebase Realtime Database
at a near-real-time rate of up to 1\,Hz.

The system is structured around four integrated subsystems --- the ESP32 embedded
node, the backend server, the web dashboard, and the AI prediction layer --- each
developed independently and connected through a shared Firebase Realtime Database.
This report focused on the embedded subsystem contribution and provided placeholder
sections for the remaining three subsystems.

% ============================================================
\section{Results Summary}
% ============================================================

The key results achieved by the embedded subsystem are:

\begin{itemize}
	\item Sensor accuracy within manufacturer specifications: $\pm$0.1\,°C for
	      temperature, $\pm$0.4\,\% for humidity, $\pm$1.0\,BPM for heart rate, and
	      $\pm$1.0\,\% for SpO\textsubscript{2}.
	\item Steady-state end-to-end latency of $\approx$2\,200\,ms for room data
	      (dominated by the DHT22 hardware sampling period) and $\approx$250\,ms for
	      patient data.
	\item Successful automatic Wi-Fi recovery in both mid-session dropout and
	      cold-start failure scenarios, with reconnection within 10--15\,s.
	\item Stable email/password Firebase authentication with full event traceability
	      via the async callback, confirmed across all test sessions.
	\item Zero invalid records written to the database across all sessions, thanks
	      to the NaN guard and IR threshold checks.
	\item Complete configuration persistence across five cold reboots via NVS.
\end{itemize}

% -------------------------------------------------------
% TODO (Backend Team): add a paragraph summarising your
% key results (API response times, alert accuracy, etc.).
% -------------------------------------------------------
\begin{center}
	\textit{[Backend results summary to be added.]}
\end{center}

% -------------------------------------------------------
% TODO (Web Dashboard Team): add a paragraph summarising
% your key results.
% -------------------------------------------------------
\begin{center}
	\textit{[Web dashboard results summary to be added.]}
\end{center}

% -------------------------------------------------------
% TODO (AI Team): add a paragraph summarising your key
% results (model accuracy, inference latency, etc.).
% -------------------------------------------------------
\begin{center}
	\textit{[AI results summary to be added.]}
\end{center}

% ============================================================
\section{Limitations}
% ============================================================

The primary limitation of the embedded subsystem is the absence of a local offline
buffer: records generated during network outages are permanently lost.
Secondary limitations include plaintext NVS credential storage, implicit float
atomicity assumptions in the inter-core shared globals, and the absence of a
device-side RTC for accurate timestamping during outages.

% -------------------------------------------------------
% TODO (all teams): add a sentence per subsystem summarising
% the main unresolved limitation.
% -------------------------------------------------------
\begin{center}
	\textit{[Limitations for all other subsystems to be summarised here.]}
\end{center}

% ============================================================
\section{Future Work}
% ============================================================

The following directions are identified for future development of the complete system:

\begin{enumerate}
	\item \textbf{Multi-room scalability} (backend + embedded): extend the data
	      model and backend to support concurrent streams from multiple ESP32 nodes
	      deployed across different hospital rooms.

	\item \textbf{Real-time push notifications} (backend + web): replace the
	      polling-based dashboard with WebSocket or Server-Sent Events for sub-second
	      alert delivery to medical staff.

	\item \textbf{Enhanced AI model} (AI): extend the prediction model to
	      incorporate multi-parameter correlation (e.g.\ combined BPM, SpO\textsubscript{2},
	      and room temperature trends) to improve anomaly detection sensitivity and
	      reduce false positive rates.

	\item \textbf{Motion artefact rejection} (embedded): add an accelerometer to
	      suppress SpO\textsubscript{2} readings during patient movement.

	\item \textbf{Power management} (embedded): implement Wi-Fi modem sleep to
	      reduce power consumption for potential battery-backup deployments.
\end{enumerate}
